\documentclass[11pt]{amsart}
\usepackage[margin=1.12in]{geometry}
\usepackage[T1]{fontenc}
\usepackage{amsmath,amssymb,amsthm}
\setlength{\emergencystretch}{2em}
\DeclareMathSizes{8}{8}{8}{8}
\DeclareMathSizes{9}{9}{8}{8}
\DeclareMathSizes{10.95}{10.95}{8}{8}
\DeclareFontFamily{OMX}{cmex}{}
\DeclareFontShape{OMX}{cmex}{m}{n}{<-> cmex10}{}
\usepackage[hidelinks]{hyperref}
\newtheorem{theorem}{Theorem}
\newtheorem{corollary}[theorem]{Corollary}
\title[Qualitative quasicircle preservation]{Qualitative quasicircle preservation implies quasiconformality}
\date{}
\begin{document}
\begin{abstract}
A homeomorphism of a closed oriented Riemann surface that sends every
quasicircle to a quasicircle is quasiconformal. In positive genus,
essential quasicircles suffice. The proof combines a sequence of
large-distortion triples into one quasicircle whose image fails bounded
turning. This answers a question of Booth, Nolte and Verberne and identifies
the full automorphism group of their essential quasicircle graph.
\end{abstract}
\maketitle

\section{Question and statement}
Booth, Nolte and Verberne~\cite[Section 1]{BNV} ask whether every
homeomorphism of a closed surface sending quasicircles to quasicircles
is quasiconformal. Their graph theorem applies in genus at least two
and permits orientation reversal. We use the same convention for
quasiconformality.

\begin{theorem}\label{main}
Let \(S\) be a closed oriented Riemann surface and let \(f:S\to S\) be
a homeomorphism. If \(f\) takes each quasicircle to a quasicircle, then
\(f\) is quasiconformal. When \(S\) has positive genus, it suffices to
test essential quasicircles.
\end{theorem}

The image quality bound may depend on the entire input curve. Only
forward preservation is assumed. Gehring's Definition~12 for plane
self-homeomorphisms~\cite{Gehring} asks for a uniform supremum of a
four-point constant over images of round circles. Aseev's
criteria~\cite{Aseev2014,Aseev2016} likewise require a common bound.
Booth, Nolte and Verberne~\cite[Section 5.3]{BNV} put each quadruple
on its own essential curve and use weak coarse order preservation
to obtain a common image bound. Our construction puts infinitely
many distortion witnesses on a single essential curve.

\section{Proof of Theorem~\ref{main}}
An embedded circle \(C\) has bounded turning if for every \(x,y\in C\)
the smaller diameter of its two subarcs from \(x\) to \(y\) is
at most \(K d(x,y)\), with \(K\) independent of \(x,y\). We use
the equivalence with quasicircularity, by the classical Ahlfors
criterion~\cite{Ahlfors} in the plane or sphere and by
Remark~1 and Lemma~2.8 of~\cite{BNV} for essential curves.
Distances in a relatively compact subdisk of a conformal chart
are comparable with Euclidean distances after shrinking the disk.
All smooth metrics on the compact surface are bilipschitz equivalent.

\paragraph{Distinct centers.}
Suppose \(f\) is not quasiconformal. Choose a relatively compact open
set \(W\) in a conformal chart whose closure is mapped into a target
chart, on which the coordinate map \(F\) is not quasiconformal.
A finite cover by such coordinate neighborhoods would give a global
bound if every restriction were quasiconformal. Define
\[
 H(z)=\limsup_{r\downarrow0}
 \frac{\max_{|u-z|=r}|F(u)-F(z)|}{\min_{|u-z|=r}|F(u)-F(z)|}.
\]
For every \(M\) there are infinitely many \(z\in W\) with \(H(z)>M\).
Otherwise the metric criterion~\cite{Vaisala} gives a uniform
quasiconformal bound on \(W\) off finitely many points. These points
are removable. In the modulus formulation the families of curves
meeting the punctures or their images have modulus zero. Near each
puncture, arbitrarily small multiples of \(1/(r\log(1/r))\) have
arbitrarily small \(L^2\) norm and infinite integral along any
nonconstant rectifiable curve reaching the puncture. Adding such
zero-modulus families back preserves the modulus comparisons.

Choose distinct \(z_j\) with \(H(z_j)\to\infty\) and pass to a
subsequence converging to \(z_*\in\overline W\), deleting a center
equal to \(z_*\) if needed. After translating \(z_*\) to \(0\),
taking a convergent subsequence of approach directions, and rotating,
write \(z_j=(t_j,y_j)\) with \(t_j>0\), \(|y_j|\le t_j/4\) and
\(t_{j+1}<t_j/4\). Choose \(r_j<t_j/100\) so that the max/min ratios
on the circles of radius \(r_j\) tend to infinity. Shrink the radii
further within the available arbitrarily small distortion witnesses
so that \(B(z_j,8r_j)\) lies in a fixed small chart rectangle about
\(0\).

\paragraph{Separated triples.}
Start with a maximum and a minimum point on each circle. If the
angle between their directions is outside \([\pi/4,3\pi/4]\), use
an intermediate point at a right angle to the maximum direction.
Its angle from the minimum direction lies in \([\pi/4,3\pi/4]\)
as well. At least one of the two successive ratios through the
intermediate point is at least the square root of the original
ratio. Thus we obtain points \(a_j,b_j\) with the stated angle
separation and
\[
 \frac{|F(a_j)-F(z_j)|}{|F(b_j)-F(z_j)|}\longrightarrow\infty.
\]

\paragraph{Uniform local patch.}
Normalize \(z=0\) and \( |a|=|b|=1\). Let \(\theta\) be the angle
between \(a\) and \(b\), and \(e\) the unit bisector of the directed
segments \(a\) and \((b-a)/|b-a|\). Their dot product is
\(-\sin(\theta/2)\), so their projections onto \(e\) are both
at least \(\sin(\pi/16)\). In coordinates with horizontal axis \(e\),
\(0\to a\to b\) is the graph of a function \(h\) of uniformly
bounded slope and height at most one, with
\(0<e\cdot a<e\cdot b\le1\). Extend \(h\) by zero to the left of \(0\),
linearly to zero at \(t=2\) on the right, and by zero thereafter.
The map
\[
 G(t,u)=(t,u+h(t)\rho(u)),\qquad
 \rho(u)=\max(0,1-|u|/2)
\]
is uniformly bilipschitz. Each vertical fiber is sent onto itself
by an increasing map with a lower Lipschitz bound \(1/2\).
The uniform Lipschitz bound in \(t\) then gives the inverse
Lipschitz estimate on pairs from different fibers. Its support
lies in \(B(0,4)\).

Let \(R\) rotate the original horizontal axis to \(e\). Let \(T\)
be the radial twist equal to \(R\) up to radius \(4\), tapering
linearly in angle to the identity at radius \(8\), and equal to
the identity thereafter. Choose the rotation angle with absolute
value at most \(\pi\). Comparing the rotating vectors first at
the smaller radius bounds the error by the smaller radius times
the angle difference. This gives a uniform Lipschitz bound, and
the same estimate with the opposite angle controls the inverse.
The patch \(R\circ G\circ R^{-1}\circ T\) is uniformly bilipschitz,
fixes the exterior of \(B(0,8)\), and sends the diameter through
\(0\) to an arc containing \(0,a,b\) in that order. Indeed, on
\(0\le t\le e\cdot b\) the image of \((t,0)\) is \(R(t,h(t))\).
Rescale and translate the patch to each \(B(z_j,8r_j)\).

\paragraph{A single witness curve.}
Choose a smooth embedded circle meeting a sufficiently small
rectangle about \(0\) in just a horizontal arc, and fix \(w\) on
the circle outside that rectangle. In positive genus start with
a smooth essential circle, move a point and its tangent by
smooth ambient isotopies, and straighten the local graph with
a supported smooth shear. A contractible circle with the same
local geometry also serves for the all-quasicircle test on the
sphere.

After discarding a finite prefix, replace the horizontal arc
with a graph whose height is \(y_j\) on each disjoint interval
\([t_j-8r_j,t_j+8r_j]\). Interpolate linearly on the bounded gaps
and to zero at the positive end, and take height zero at and to
the left of \(0\). The bounds on \(t_j,y_j,r_j\) give uniformly
bounded slopes on the gaps and continuity at \(0\), so the
graph is Lipschitz across \(0\) as well. For example the gap
between consecutive plateaus is greater than \(0.65t_j\) and
the height difference is at most \(5t_j/16\). Retaining sufficiently
small centers keeps the graph away from the top and bottom edges.
Multiplying the height by a parameter gives an isotopy relative
to the arc endpoints. The resulting circle is a quasicircle by
bounded turning on local Lipschitz graphs and compactness, and it
crosses each support disk along the horizontal diameter.

Apply the patches on the pairwise disjoint disks and the identity
elsewhere. The pasted map is bilipschitz. For points in different
disks, split their straight segment at the disks' fixed boundaries.
Summing the estimates on those segments gives the common Lipschitz
bound. For one point outside the disks one split suffices, and the
same proof for the inverse gives the lower bound. Each disk is
preserved and the family accumulates inside the chart. The map
extends by identity to the surface and preserves the isotopy
class, since in each disk the supported homeomorphism is isotopic
to the identity relative to its boundary by radial contraction,
and the support diameters tend to zero. The image curve \(C\)
is a quasicircle containing \(z_j,a_j,b_j\) in order on each
diameter replacement. The complementary arc from \(b_j\) to
\(z_j\) contains \(w\).

Now the \(f(a_j)\)-containing arc from \(f(z_j)\) to \(f(b_j)\)
has diameter at least a fixed positive multiple of
\(|F(a_j)-F(z_j)|\). The complementary arc contains \(f(w)\)
and its diameter is bounded below by a positive constant since
\(f(z_j)\to f(z_*)\ne f(w)\). The endpoint distance is comparable
with \(|F(b_j)-F(z_j)|\), which tends to zero. The diameter of each
arc divided by the endpoint distance therefore tends to infinity.
This contradicts bounded turning of \(f(C)\) and proves the theorem.

\section{Worked example}
Consider the polygon with vertices \(z=(0,0),a=(1,0),b=(0,1)\),
and \(w=(-1,-1)\) in that cyclic order. For \(L\ge1\) let
\(A_L(x,y)=(Lx,y)\). The two arcs of the image from \(A_L(z)\)
to \(A_L(b)\) contain \(A_L(a)\) and \(A_L(w)\), respectively,
so their diameters are at least \(L\) and \(\sqrt{L^2+1}\).
Their common endpoint distance is one. At \(L=8\) this pair forces
a bounded-turning constant of at least \(8\). Each fixed \(A_L\)
is quasiconformal. The example explains the turning estimate for
one triple, while the proof produces arbitrarily large witnesses
for a single \(f\) and a single \(C\).

\begin{corollary}
For a closed oriented surface \(S\) of genus at least two,
\(\mathcal Q(S)\to\operatorname{Aut}\mathcal{QC}^{\dagger}(S)\)
is an isomorphism.
\end{corollary}
\begin{proof}
The first assertion of Theorem A in~\cite{BNV} represents every
graph automorphism by a homeomorphism. It takes every essential
quasicircle to a quasicircle, so Theorem~\ref{main} makes the
representing map quasiconformal. Faithfulness follows from
Theorem A, which also implies that every graph automorphism
is coarsely order preserving.
\end{proof}

\bibliographystyle{amsplain}
\begingroup\raggedright
\bibliography{references}
\endgroup
\end{document}
